% ============================================================================
%  CHAPITRE 5 -- PROTOCOLE D'ÉVALUATION ET MÉTRIQUES
% ============================================================================
\chapter{Protocole d'évaluation et métriques}
\label{chap:protocole}

Mesurer « le sens » est le point délicat : un protocole mal construit produirait des chiffres flatteurs et
faux.

\section{Jeux de données}
\label{sec:donnees}

\begin{table}[H]
\centering
\caption{Jeux de données et rôle de chacun}
\label{tab:donnees}
\small
\begin{tabularx}{\textwidth}{@{}p{2.1cm}p{2.4cm}p{2.4cm}>{\raggedright\arraybackslash}X@{}}
\toprule
\textbf{Jeu} & \textbf{Taille} & \textbf{Usage} & \textbf{Remarque} \\
\midrule
CIFAR-10~\cite{krizhevsky2009} & $32\times32$, 10 classes & entraînement et évaluation & \num{50000} images d'apprentissage ; test sur les \num{5000} premières, pas de \SI{1}{\decibel} \\
Kodak~\cite{kodak} & $768\times512$, 24 images & évaluation seule & jamais vues à l'entraînement ; pas de \SI{2}{\decibel} \\
Vidéo CIF~\cite{xiph} & $352\times288$, 3 clips & évaluation seule & 30 premières trames d'\srcfile{akiyo} (statique), \srcfile{foreman} (modéré), \srcfile{bus} (rapide) \\
\bottomrule
\end{tabularx}
\end{table}

Les clips CIF remplacent les séquences UVG~\cite{uvg} suggérées initialement, dont les fichiers 1080p bruts
pèsent plusieurs gigaoctets. Les trois régimes de mouvement retenus s'avèrent déterminants pour l'extension
vidéo du chapitre~\ref{chap:extensions}.

\section{Métriques}
\label{sec:metriques}

\paragraph{Fidélité.} Le PSNR, en décibels,
\begin{equation}
  \mathrm{PSNR}(\hat{x},x) = 10\log_{10}\frac{1}{\mathrm{MSE}(\hat{x},x)} \quad \text{(images dans } [0,1]),
  \label{eq:psnr}
\end{equation}
vaut \SI{30}{\decibel} pour une bonne reconstruction et \SI{20}{\decibel} pour une image visiblement
dégradée. Le SSIM~\cite{wang2004}, entre 0 et 1, compare localement luminance, contraste et structure.

\paragraph{Sens.} Sur CIFAR-10, la \textbf{précision} d'un classifieur sur l'image reçue ; la référence sur
images originales est \textbf{88,4\,\%}. Sur Kodak et en vidéo, faute d'étiquette, un
ResNet-18~\cite{he2016} pré-entraîné sur ImageNet fournit la similarité cosinus des caractéristiques
(\srcfile{sem_cos}) et l'accord des étiquettes avant/après transmission (\srcfile{top1_agree}).

\section{L'indépendance des juges}
\label{sec:juges}

\begin{keybox}[Règle de séparation]
Le réseau qui intervient dans la perte n'est jamais celui qui mesure le résultat.
\end{keybox}

Sans cette séparation, le système sémantique pourrait apprendre à produire des artefacts qui plaisent à son
propre correcteur : la précision mesurée monterait sans que le sens soit mieux transmis.

\begin{table}[H]
\centering
\caption{Rôle de chaque réseau : apprentissage ou mesure, jamais les deux}
\label{tab:juges}
\small
\begin{tabularx}{\textwidth}{@{}p{3.4cm}ccX@{}}
\toprule
\textbf{Réseau} & \makecell{\textbf{Dans une}\\\textbf{perte}} & \makecell{\textbf{Utilisé pour}\\\textbf{mesurer}} & \textbf{Précision sur images originales} \\
\midrule
Classifieur A (ResNet) & \cmark & \xmark & 85,4\,\% \\
Classifieur B (VGG) & \xmark & \cmark & 88,4\,\% \\
ResNet-18 ImageNet & \xmark & \cmark & référence de similarité sémantique \\
\bottomrule
\end{tabularx}
\end{table}

\section{Conduite des mesures}
\label{sec:mesures}

Pour chaque couple (MCS, SNR), les paquets de toutes les images sont codés, modulés, envoyés et décodés par
Viterbi. Les graines sont fixées par point de mesure. Chaque scheme produit une ligne de CSV dans
\srcfile{results/eval/} ; tableaux et figures en sont dérivés par \srcfile{scripts/make_figures.py}, sans
recopie manuelle.
